\section{Technology Stack and Integration}
\label{sec:tech}

Every part of \proj{} runs on the client, inside the browser. The system is
assembled from a small set of web technologies, each chosen for one job, and the
engineering value is largely in how they are \emph{coordinated}. This section
catalogues the stack, explains how each piece is used, and then shows how they
fit together.

\subsection{The Stack at a Glance}
Table~\ref{tab:stack} lists each technology, the single job it performs, and the
file where it is wired into the system; the subsections that follow expand on how
each one is used.

\begin{table}[ht]
  \centering
  \caption{Technologies used in \proj{}, the single role each plays, and the source
           file where it is wired into the system.}
  \label{tab:stack}
  \renewcommand{\arraystretch}{1.3}
  {\small
  \begin{tabularx}{\textwidth}{@{}>{\raggedright\arraybackslash}p{3.7cm} X >{\raggedright\arraybackslash}p{2.6cm}@{}}
    \toprule
    \textbf{Technology} & \textbf{Role in the project} & \textbf{Where} \\
    \midrule
    Chrome Extension (MV3) & Packaging, permissions, and the CSP that lets WASM OCR run & \code{manifest.json} \\
    Service worker + \code{chrome.\allowbreak action} / \code{chrome.\allowbreak commands} & Trigger, per-tab progress badge, tab switch, and serving the searchable copy & \code{service-\allowbreak worker.js} \\
    Offscreen document & Hidden page with a DOM that hosts the whole pipeline & \code{offscreen/\allowbreak offscreen.js} \\
    PDF.js & Detects text and images per page; renders pages for OCR & \code{processor.js}, \code{ocr.js} \\
    HTML5 Canvas & OCR raster target and pixel-level preprocessing & \code{ocr.js} \\
    Tesseract.js (WASM) + Web Workers & Optical character recognition on a pool of workers & \code{ocr.js} \\
    pdf-lib & Writes the invisible text layer into the copy and saves it & \code{embed.js} \\
    Web Crypto (SHA-256) & Content-based document identity & \code{cache.js} \\
    IndexedDB & Persistent OCR result cache with LRU pruning & \code{cache.js} \\
    Cache Storage & Holds searchable copies, served at extension URLs & \code{offscreen.js}, \code{service-\allowbreak worker.js} \\
    \code{chrome.storage}, \code{chrome.\allowbreak contextMenus} & OCR language and open copies; the language menu & \code{service-\allowbreak worker.js} \\
    Chrome's native PDF viewer & All display, find, highlighting, selection, and printing & (Chrome) \\
    \bottomrule
  \end{tabularx}}
\end{table}

\subsection{How Each Technology Is Used}
\begin{description}[leftmargin=1.4em, style=nextline]
  \item[Chrome Extension (Manifest V3).] \code{manifest.json} requests the
        \code{offscreen}, \code{storage} and \code{contextMenus} permissions and
        host access to fetch PDFs, sets a content security policy of
        \code{script-src 'self' 'wasm-\allowbreak unsafe-\allowbreak eval'} so
        Tesseract's WebAssembly core can execute, and requires Chrome~116 or
        later (for \code{chrome.\allowbreak runtime.\allowbreak getContexts}). There are no
        content scripts and no web-accessible resources: the extension never
        touches web pages.
  \item[Service worker, \code{chrome.action} \& \code{chrome.commands}.] The
        service worker sleeps until the user clicks the toolbar icon
        (\code{action.onClicked}) or presses the shortcut
        (\code{commands.onCommand}, \code{Ctrl+Shift+F}; Control+Shift+F on Mac).
        It starts the offscreen document, sends it the job, mirrors progress as a
        per-tab badge (an ellipsis, then a percentage, then a check mark, or a red \code{!} whose
        tooltip explains the error), and finally navigates the tab to the copy.
        Its \code{fetch} handler serves the copy's bytes, and invoking the
        extension on a copy restores the original URL.
  \item[Offscreen document.] Created with reason \code{WORKERS}, it provides the
        DOM, canvas and Web Workers the pipeline needs. It processes one PDF at a
        time from a queue and releases its OCR workers after a minute idle.
  \item[PDF.js.] \code{getDocument(\{data\})} parses the downloaded bytes. Per
        page, \code{getTextContent()} returns existing text,
        \code{getOperatorList()} (checked against the \code{pdfjsLib.OPS} image
        operators) reveals raster images, and \code{render()} rasterizes pages
        that need OCR. \code{viewport.\allowbreak convertToPdfPoint} maps OCR pixels
        back into PDF user space. PDF.js never displays anything to the user.
  \item[Canvas API.] Each page needing OCR is rendered to an offscreen canvas
        about 3000\,px wide, preprocessed through \code{getImageData} /
        \code{putImageData}, and released immediately after recognition.
  \item[Tesseract.js.] Loaded as a UMD script exposing \code{Tesseract}. A
        \code{createScheduler()} fronts a pool of LSTM \code{createWorker()}
        workers configured with \code{PSM.SPARSE\_TEXT}; \code{addJob("recognize")}
        with \code{blocks:true} returns the block\,$\to$\,word hierarchy with
        bounding boxes. The WASM core and the \code{eng}/\code{dan} language data
        are bundled, so OCR is fully offline. Only the SIMD and plain LSTM cores
        are shipped; the relaxed-SIMD core is disabled (Section~\ref{sec:process}).
  \item[pdf-lib.] Loads the same bytes, embeds the standard Helvetica font once,
        appends one invisible text object per recognized word to each page's
        content stream, and saves the copy with metadata unchanged.
  \item[Web Crypto, IndexedDB, Cache Storage.] \code{crypto.subtle.digest}
        computes the SHA-256 of the downloaded file; IndexedDB stores each page's
        recognized words under \code{hash::lang::page}; Cache Storage holds the
        finished copy as an HTTP \code{Response}, keyed by its URL path.
  \item[\code{chrome.storage} and \code{chrome.contextMenus}.] The OCR language
        lives in \code{storage.local} and is chosen from radio items on the toolbar
        icon's right-click menu; \code{storage.session} records which tab shows
        which copy, so the copy can be deleted when that tab closes or navigates.
\end{description}

\subsection{Coordination Between Technologies}
Figure~\ref{fig:integration} traces the data that crosses each boundary. The
service worker is the coordinator between the browser and the extension; the
offscreen document's \code{processor.js} is the coordinator of the libraries.
The two exchange only small messages---a job description going in, progress and
a copy URL coming out---while the heavy bytes travel through Cache Storage, which
both can reach because they share the extension's origin.

\begin{figure}[ht]
  \centering
  \begin{tikzpicture}[
      font=\small,
      tech/.style={rectangle, rounded corners=2pt, draw=accent!70, fill=accent!6,
                   align=center, minimum height=9mm, inner sep=4pt},
      hub/.style={rectangle, rounded corners=2pt, draw=accent!90, fill=accent!12,
                  align=center, minimum height=10mm, inner sep=5pt, font=\small\bfseries},
      io/.style={rectangle, draw=rulegray, fill=codebg, align=center,
                 minimum height=9mm, inner sep=4pt},
      a/.style={-{Stealth[length=2.2mm]}, draw=accent!80, semithick},
      b/.style={{Stealth[length=2.2mm]}-{Stealth[length=2.2mm]}, draw=accent!80, semithick},
      lbl/.style={font=\scriptsize\itshape, fill=white, inner sep=1pt, align=center},
    ]
    \node[io]   (tab)    at (-5.4, 0)   {PDF tab\\(Chrome's viewer)};
    \node[hub]  (sw)     at (0, 0)      {Service worker};
    \node[tech] (cs)     at (5.2, 0)    {Cache Storage\\(copies)};
    \node[hub]  (off)    at (0, -3.0)   {Offscreen document\\\code{processor.js}};
    \node[tech] (pdfjs)  at (-5.1, -5.5) {PDF.js};
    \node[tech] (tess)   at (-1.75, -5.5) {Tesseract.js\\(Web Workers)};
    \node[tech] (pdflib) at (1.75, -5.5) {pdf-lib};
    \node[tech] (idb)    at (5.1, -5.5) {IndexedDB\\(OCR cache)};

    \draw[a] ([yshift=2mm]tab.east) -- node[lbl,above]{\code{Ctrl+Shift+F} / click} ([yshift=2mm]sw.west);
    \draw[a] ([yshift=-2mm]sw.west) -- node[lbl,below]{open copy URL} ([yshift=-2mm]tab.east);
    \draw[b] (sw) -- node[lbl,above]{\code{fetch} / copy bytes} (cs);
    \draw[a] ([xshift=-4mm]sw.south) -- node[lbl,left]{job: URL,\\language, cache key} ([xshift=-4mm]off.north);
    \draw[a] ([xshift=4mm]off.north) -- node[lbl,right]{progress,\\done / error} ([xshift=4mm]sw.south);
    \draw[a] (off.east) -- node[lbl,right=2pt]{store copy} (cs.south west);
    \draw[b] (off) -- node[lbl,pos=0.72,above left]{detect / render} (pdfjs);
    \draw[b] (off) -- node[lbl,left]{canvas /\\words + boxes} (tess);
    \draw[a] (off) -- node[lbl,right]{invisible\\words} (pdflib);
    \draw[b] (off) -- node[lbl,above right]{cached words} (idb);
  \end{tikzpicture}
  \caption{Integration view. The service worker talks to the browser (trigger,
           tab navigation, serving the copy); the offscreen document drives the
           libraries. Arrow labels name the data crossing each boundary;
           double-headed arrows indicate request--response pairs. The finished
           copy reaches the viewer through Cache Storage, not through messages.}
  \label{fig:integration}
\end{figure}
